\section{AIME I 1997}
        \begin{enumerate}
        	\item (\textit{AIME I 1997}) How many of the integers between 1 and 1000, inclusive, can be expressed as the difference of the squares of two nonnegative integers?


 
 \textbf{Answer: }750



\textbf{Solution}

Notice that all odd numbers can be obtained by using $(a+1)^2-a^2=2a+1,$ where $a$ is a nonnegative integer. All multiples of $4$ can be obtained by using $(b+1)^2-(b-1)^2 = 4b$, where $b$ is a positive integer. Numbers congruent to $2 \pmod 4$ cannot be obtained because squares are $0, 1 \pmod 4.$ Thus, the answer is $500+250 = \boxed{750}.$



	\item (\textit{AIME I 1997}) The nine horizontal and nine vertical lines on an $8\times8$ checkerboard form $r$ rectangles, of which $s$ are squares.  The number $s/r$ can be written in the form $m/n,$ where $m$ and $n$ are relatively prime positive integers.  Find $m + n.$

 \textbf{Answer: }125



\textbf{Solution}

To determine the two horizontal sides of a rectangle, we have to pick two of the horizontal lines of the checkerboard, or ${9\choose 2} = 36$. Similarly, there are ${9\choose 2}$ ways to pick the vertical sides, giving us $r = 1296$ rectangles.

 For $s$, there are $8^2$ \href{https://artofproblemsolving.com/wiki/index.php?title=Unit\_square}{unit squares}, $7^2$ of the $2\times2$ squares, and so on until $1^2$ of the $8\times 8$ squares. Using the sum of squares formula, that gives us $s=1^2+2^2+\cdots+8^2=\dfrac{(8)(8+1)(2\cdot8+1)}{6}=12*17=204$.

 Thus $\frac sr = \dfrac{204}{1296}=\dfrac{17}{108}$, and $m+n=\boxed{125}$.

 



\textbf{Solution 2}

First, to find the number of squares, we can look case by case by the side length of the possible squares on the checkerboard. We see that there are $8^2$ ways to place a $1$ x $1$ square and $7^2$ for a $2$ x $2$ square. This pattern can be easily generalized and we see that the number of squares is just $\sum^8_{i=1}{i^2}$. This can be simplified by using the well-known formula for the sum of consecutive squares $\frac{n(n+1)(2n+1)}{6}$ to get $204$.

 Then, to find the number of rectangles, first note that a square falls under the definition of a rectangle. We can break up the rectangles into cases for the length x width. As we note down the cases for $1$x$1, 1$x$2 , 2$x$1, 2$x$2,...,$ we see they are respectively $8$x$8, 8$x$7, 7$x$8, 7$x$7, ...$. We can quickly generalize this pattern to basically just ${\sum^8_{i=1}{i}}\cdot{\sum^8_{i=1}{i}}$. This gets us ${(\frac{9\cdot8}{2})}^2,$ which is just $1296.$

 Now, to calculate the ratio of $s/r,$ we divide $204$ by $1296$ to get a simplified fraction of $\frac{17}{108}.$

 Thus, our answer is just $s + r = 17+108 = \boxed{125}$~MathWhiz35

 


	\item (\textit{AIME I 1997}) Sarah intended to multiply a two-digit number and a three-digit number, but she left out the multiplication sign and simply placed the two-digit number to the left of the three-digit number, thereby forming a five-digit number.  This number is exactly nine times the product Sarah should have obtained.  What is the sum of the two-digit number and the three-digit number?


 
 \textbf{Answer: }126



\textbf{Solution}

Let $x$ be the two-digit number, $y$ be the three-digit number. Putting together the given, we have $1000x+y=9xy \Longrightarrow 9xy-1000x-y=0$. Using \href{https://artofproblemsolving.com/wiki/index.php?title=Simon\%27s\_Favorite\_Factoring\_Trick}{SFFT}, this factorizes to $(9x-1)\left(y-\dfrac{1000}{9}\right)=\dfrac{1000}{9}$, and $(9x-1)(9y-1000)=1000$.

 Since $89 < 9x-1 < 890$, we can use trial and error on factors of 1000. If $9x - 1 = 100$, we get a non-integer. If $9x - 1 = 125$, we get $x=14$ and $y=112$, which satisifies the conditions. Hence the answer is $112 + 14 = \boxed{126}$.

 



\textbf{Solution 2}

As shown above, we have $1000x+y=9xy$, so $1000/y=9-1/x$. $1000/y$ must be just a little bit smaller than 9, so we find $y=112$, $x=14$, and the solution is $\boxed{126}$.

 


 



\textbf{Solution 3}

To begin, we rewrite $(10a+b)*(100x+10y+z)*9 = 10000a + 1000b + 100x + 10y + z$

 as

 $(90a+9b-1)(100x+10y+z) = 10000a + 1000b$

 and

 $(90a+9b-1)(100x+10y+z) = 1000(10a + b)$

 
This is the most important part: Notice $(90a+9b-1)$ is $-1 \pmod{10a+b}$ and $1000(10a + b)$ is $0\pmod{10a+b}$. That means $(100x+10y+z)$ is also $0\pmod{10+b}$. Rewrite $(100x+10y+z)$ as $n\times(10a+b)$.

 
$(90a+9b-1)\times n(10a+b)= 1000(10a + b)$
$(90a+9b-1)\times n= 1000$

 
Now we have to find a number that divides 1000 using prime factors 2 or 5 and is $8\pmod9$. It is quick to find there is only one: 125. That gives 14 as $10a+b$ and 112 as $100x+10y+z$. Therefore the answer is $112 + 14 = \boxed{126}$.

 


 -jackshi2006

 


	\item (\textit{AIME I 1997}) Circles of radii 5, 5, 8, and $m/n$ are mutually externally tangent, where $m$ and $n$ are relatively prime positive integers.  Find $m + n.$

 \textbf{Answer: }017



\textbf{Solution}


                    \begin{figure}[H]
                        \centering
                        \includegraphics[width=0.5\textwidth]{Solution/AIME_I_1997_4_Solution_Figure_1.png}
                    \end{figure}
                    

 If (in the diagram above) we draw the line going through the centers of the circles with radii $8$ and $\frac mn = r$, that line is the perpendicular bisector of the segment connecting the centers of the two circles with radii $5$. Then we form two \href{https://artofproblemsolving.com/wiki/index.php?title=Right\_triangles}{right triangles}, of lengths $5, x, 5+r$ and $5, 8+r+x, 13$, wher $x$ is the distance between the center of the circle in question and the segment connecting the centers of the two circles of radii $5$. By the \href{https://artofproblemsolving.com/wiki/index.php?title=Pythagorean\_Theorem}{Pythagorean Theorem}, we now have two equations with two unknowns:

 
 \begin{eqnarray*} 5^2 + x^2 &=& (5+r)^2 \\ x &=& \sqrt{10r + r^2} \\ && \\ (8 + r + \sqrt{10r+r^2})^2 + 5^2 &=& 13^2\\ 8 + r + \sqrt{10r+r^2} &=& 12\\ \sqrt{10r+r^2}&=& 4-r\\ 10r+r^2 &=& 16 - 8r + r^2\\ r &=& \frac{8}{9} \end{eqnarray*}

 So $m+n = \boxed{17}$.

 
NOTE: It can be seen that there is no apparent need to use the variable x as a 5,12,13 right triangle has been formed.

 Using the 5,12,13 triangle will lead to the equation $5^2+(12-8-r)^2=(5+r)^2$ which yields the same solution.

 



\textbf{Solution 2}

We may also use Descartes' theorem, $k_4=k_1+k_2+k_3\pm 2\sqrt{k_1k_2+k_2k_3+k_3k_1}$ where each of $k_i$ is the curvature of a circle with radius $r_i$, and the curvature is defined as $k_i=\frac{1}{r_i}$. The larger solution for $k_4$ will give the curvature of the circle externally tangent to the other circles, while the smaller solution will give the curvature for the circle internally tangent to each of the other circles. Using Descartes' theorem, we get $k_4=\frac15+\frac15+\frac18+2\sqrt{\frac{1}{40}+\frac{1}{40}+\frac{1}{25}}=\frac{21}{40}+2\sqrt{\frac{45}{500}}=\frac{45}{40}$. Thus, $r_4=\frac{1}{k_4}=\frac{40}{45}=\frac89$, and the answer is $\boxed{017}$



	\item (\textit{AIME I 1997}) The number $r$ can be expressed as a four-place decimal $0.abcd,$ where $a, b, c,$ and $d$ represent digits, any of which could be zero.  It is desired to approximate $r$ by a fraction whose numerator is 1 or 2 and whose denominator is an integer. The closest such fraction to $r$ is $\frac 27.$  What is the number of possible values for $r$?


 
 \textbf{Answer: }417



\textbf{Solution}

The nearest fractions to $\frac 27$ with numerator $1$ are $\frac 13, \frac 14$; and with numerator $2$ are $\frac 26, \frac 28 = \frac 13, \frac 14$ anyway. For $\frac 27$ to be the best approximation for $r$, the decimal must be closer to $\frac 27 \approx .28571$ than to $\frac 13 \approx .33333$ or $\frac 14 = .25$.

 Thus $r$ can range between $\frac{\frac 14 + \frac{2}{7}}{2} \approx .267857$ and $\frac{\frac 13 + \frac{2}{7}}{2} \approx .309523$. At $r = .2678, .3096$, it becomes closer to the other fractions, so $.2679 \le r \le .3095$ and the number of values of $r$ is $3095 - 2679 + 1 = \boxed{417}$.

 


	\item (\textit{AIME I 1997}) Point $B$ is in the exterior of the regular $n$-sided polygon $A_1A_2\cdots A_n$, and $A_1A_2B$ is an equilateral triangle. What is the largest value of $n$ for which $A_1$, $A_n$, and $B$ are consecutive vertices of a regular polygon?


 
 \textbf{Answer: }042



\textbf{Solution 1}


                    \begin{figure}[H]
                        \centering
                        \includegraphics[width=0.5\textwidth]{Solution/AIME_I_1997_6_Solution_Figure_1.png}
                    \end{figure}
                    

 Let the other regular polygon have $m$ sides. Using the interior angle of a regular polygon formula, we have $\angle A_2A_1A_n = \frac{(n-2)180}{n}$, $\angle A_nA_1B = \frac{(m-2)180}{m}$, and $\angle A_2A_1B = 60^{\circ}$. Since those three angles add up to $360^{\circ}$,

 
 \begin{eqnarray*} \frac{(n-2)180}{n} + \frac{(m-2)180}{m} &=& 300\\ m(n-2)180 + n(m-2)180 &=& 300mn\\ 360mn - 360m - 360n &=& 300mn\\ mn - 6m - 6n &=& 0 \end{eqnarray*}Using \href{https://artofproblemsolving.com/wiki/index.php?title=Simon\%27s\_Favorite\_Factoring\_Trick}{SFFT}, 

 
 \begin{eqnarray*} (m-6)(n-6) &=& 36 \end{eqnarray*}Clearly $n$ is maximized when $m = 7, n = \boxed{042}$.

 



\textbf{Solution 2}

As above, find that $mn - 6m - 6n = 0$ using the formula for the interior angle of a polygon.  

 Solve for $n$ to find that $n = \frac{6m}{m-6}$.  Clearly, $m>6$ for $n$ to be positive.

 With this restriction of $m>6$, the larger $m$ gets, the smaller the fraction $\frac{6m}{m-6}$ becomes.  This can be proven either by calculus, by noting that $n = \frac{6m}{m-6}$ is a  transformed hyperbola, or by dividing out the rational function to get $n = 6 + \frac{36}{m - 6}.$

 Either way, minimizng $m$ will maximize $n$, and the smallest integer $m$ such that $n$ is positive is $m=7$, giving $n = \boxed{042}$




\textbf{Solution 3}

From the formula for the measure for an individual angle of a regular n-gon, $180 - \frac{360}{n}$, the measure of $\angle A_2A_1A_n = 180 - \frac{360}{n}$. Together with the fact that an equilateral triangle has angles measuring 60 degrees, the measure of $\angle A_nA_1B = 120 + \frac{360}{n}$ (Notice that this value decreases as $n$ increases; hence, we are looking for the least possible value of $\angle A_nA_1B$). For $A_n, A_1, B$ to be vertices of a regular polygon, $\angle A_nA_1B$ must be of the form $180 - \frac{360}{n}$, where $n$ is a natural number greater than or equal to 3. It is obvious that $\angle A_nA_1B > 120$. The least angle satisfying this condition is $180 - \frac{360}{7}$. Equating this with $120 + \frac{360}{n}$ and solving yields $n = \boxed{042}$



	\item (\textit{AIME I 1997}) A car travels due east at $\frac 23$ miles per minute on a long, straight road. At the same time, a circular storm, whose radius is $51$ miles, moves southeast at $\frac 12\sqrt{2}$ miles per minute. At time $t=0$, the center of the storm is $110$ miles due north of the car. At time $t=t_1$ minutes, the car enters the storm circle, and at time $t=t_2$ minutes, the car leaves the storm circle. Find $\frac 12(t_1+t_2)$.


 
 \textbf{Answer: }198



\textbf{Solution 1}

We set up a coordinate system, with the starting point of the car at the \href{https://artofproblemsolving.com/wiki/index.php?title=Origin}{origin}. At time $t$, the car is at $\left(\frac 23t,0\right)$ and the center of the storm is at $\left(\frac{t}{2}, 110 - \frac{t}{2}\right)$. Using the distance formula, 

 
 \begin{eqnarray*} \sqrt{\left(\frac{2}{3}t - \frac 12t\right)^2 + \left(110-\frac{t}{2}\right)^2} &\le& 51\\ \frac{t^2}{36} + \frac{t^2}{4} - 110t + 110^2 &\le& 51^2\\ \frac{5}{18}t^2 - 110t + 110^2 - 51^2 &\le& 0\\ \end{eqnarray*}

 Noting that $\frac 12 \left(t_1+t_2 \right)$ is at the maximum point of the parabola, we can use $-\frac{b}{2a} = \frac{110}{2 \cdot \frac{5}{18}} = \boxed{198}$.

 



\textbf{Solution 2 (more formal explanations for Solution 1)}

First do the same process for assigning coordinates to the car. The car moves $\frac{2}{3}$ miles per minute to the right, so the position starting from $(0,0)$ is $\left(\frac{2}{3}t, 0\right)$. 

 Take the storm as circle. Given southeast movement, split the vector into component, getting position $\left(\frac{1}{2}t, 110 - \frac{1}{2}t\right)$ for the storm's center. This circle with radius 51 yields $\left(x - \frac{1}{2}t\right)^2 + \left(y -110 + \frac{1}{2}t\right)^2 = 51^2$. 

 Now substitute the car's coordinates into the circle's:

 $\left(\frac{2}{3}t - \frac{1}{2}t\right)^2 + \left(-110 + \frac{1}{2}t\right)^2 = 51^2$. 

 Simplifying and then squaring:

 $\left(\frac{1}{6}t\right)^2 + \left(-110 + \frac{1}{2}t\right)^2 = 51^2$. 

 $\frac{1}{36}t^2 + \frac{1}{4}t^2 - 110t + 110^2$

 
Forming into a quadratic we get the following, then set equal to 0, since the first time the car hits the circumference of the storm is $t_{1}$ and the second is $t_{2}$.

 
$\frac{5}{18}t^2 - 110t + 110^2 - 51^2 = 0$

 
The problem asks for sum of solutions divided by 2 so sum is equal to:

 $-\frac{b}{a} = -\frac{-110}{\frac{5}{18}} = 110\cdot{\frac{18}{5}} = 396\cdot{\frac{1}{2}} = \boxed{198}$




\textbf{Solution 3 (No Coordinates)}

We only need to know how the storm and car move relative to each other, so we can find this by subtracting the storm's movement vector from the car's.  This gives the car's movement vector as $\left(\frac{1}{6}, \frac{1}{2}\right)$.  Labeling the car's starting position A, the storm center B, and the right triangle formed by AB with a right angle at B and the car's path, we get the following diagram, with AD as our desired length since D is the average of the points where the car enters and exits the storm.

 
                    \begin{figure}[H]
                        \centering
                        \includegraphics[width=0.5\textwidth]{Solution/AIME_I_1997_7_Solution_Figure_1.png}
                    \end{figure}
                    
$AB = 110$, so $CB = \frac{110}{3}$.  The Pythagorean Theorem then gives $AC = \frac{110\sqrt{10}}{3}$, and since $\bigtriangleup ABC \sim \bigtriangleup ADB$, $AD = (AB)\frac{AB}{AC} = 33\sqrt{10}$.  The Pythagorean Theorem now gives the car's speed as $\sqrt{\frac{5}{18}}$, and finally $\frac{33\sqrt{10}}{\sqrt{\frac{5}{18}}} = \boxed{198}$.

 


	\item (\textit{AIME I 1997}) How many different $4\times 4$ arrays whose entries are all 1's and -1's have the property that the sum of the entries in each row is 0 and the sum of the entries in each column is 0?


 
 \textbf{Answer: }090



\textbf{Solution 1}


            \begin{enumerate}[label=(\roman*)]
                \item <i>For more detailed explanations, see related \href{https://artofproblemsolving.com/wiki/index.php?title=2007\_AIME\_I\_Problems/Problem\_10}{problem (AIME I 2007, 10)}.</i>
            \end{enumerate}
            The problem is asking us for all configurations of $4\times 4$ grids with 2 1's and 2 -1's in each row and column. We do casework upon the first two columns:

 
            \begin{itemize}
        \item The first two columns share no two numbers in the same row. There are ${4\choose2} = 6$ ways to pick two 1's in the first column, and the second column is determined. For the third and fourth columns, no two numbers can be in the same row (to make the sum of each row 0), so again there are ${4\choose 2}$ ways. This gives $6^2 = 36$.
        \item The first two columns share one number in the same row. There are ${4\choose 1} = 4$ ways to pick the position of the shared 1, then ${3\choose 2} = 3$ ways to pick the locations for the next two 1s, and then $2$ ways to orient the 1s. For the third and fourth columns, the two rows with shared 1s or -1s are fixed, so the only things that can be changed is the orientation of the mixed rows, in $2$ ways. This gives $4 \cdot 3 \cdot 2 \cdot 2 = 48$.
        \item The first two columns share two numbers in the same row. There are ${4\choose 2} = 6$ ways to pick the position of the shared 1s. Everything is then fixed.
            \end{itemize}
            Adding these cases up, we get $36 + 48 + 6 = \boxed{090}$.

 



\textbf{Solution 2}

Each row and column must have 2 1's and 2 -1's. Let's consider the first column. There are a total of $6$ ways to arrange 2 1's and 2 -1's. Let's consider the setup where the first and second indices of column 1 are 1 and the third and fourth are -1. Okay, now on the first row, there are 3 ways to arrange the one 1 and 2 -1's we have left to put. Now, we take cases on the second row's remaining elements. If the second row goes like 1,-1,1,-1, then by observation, there are 2 ways to complete the grid. If it goes like 1,1, -1, -1, there is 1 way to complete the grid. If it goes like 1, -1, -1, 1, then there are 2 ways to complete the grid. So our answer is $6*3*(2+1+2)$ = $\boxed{090}$.

 -pi\_is\_3.141

 



\textbf{Solution 3}

Notice that for every arrangement $A$ of the first rows of $-1$s and $1$s, we have the inverse of that row $A^{-1}$ so that the sum of the rows and columns of $A$ and $A^{-1}$ is $0$. Therefore if we have another arrangement $B$, we have $B^{-1}$. For instance, if $A=(-1,1,1,-1)$, $A^{-1}=(1,-1,-1,1)$. We then have that if we fix the first row $A$, we have that first there are $\binom{4}{2}$ values of the fixed $A$. We then have the following cases:

 Case $1$: ($AA^{-1}BB^{-1}$). $\binom{4}{2}$.

 Case $2$: ($ABA^{-1}B^{-1}$). $\binom{4}{2}$.

 Case $3$: ($AAA^{-1}A^{-1}$). $\binom{4}{2}/2$, where here we divided by $2$ because then we would overcount $AA$ and $A^{-1}A^{-1}$.

 Therefore the answer is $\binom{4}{2}(\binom{4}{2}+\binom{4}{2}+\binom{4}{2}/2)$=$6(6+6+3)$=$\boxed{090}$.

 



\textbf{Solution 4 (Bijection and Casework)}

How many orderings are possible for the first row? Well, that's easy! There are \( 4 \choose 2 \) = 6 possible ways to distribute two 1's into the first row, and the remaining -1's are uniquely determined.

 Now, the hard part. We know that to obtain a commutative sum of 0, each row and column must have two 1's and two -1's. So, we can consider the bottom 3x2 matrix.

 There are a total of \( 3 \choose 2 \) * \( 3 \choose 2 \) = 9 different orderings that satisfy the conditions of the sum to 0 (that is, \( 3 \choose 2 \) ways to distribute the -1's in the first column, and \( 3 \choose 2 \) ways to distribute the 1's in the second column, where the last number is uniquely determined). We have a couple of cases.

 Case 1: In the bottom left 3x2 matrix of the 4x4 grid, each row contains 1, -1 or -1, 1.

 This is obviously not possible as each row contains a maximum of two -1's or two 1's and by simple logic or pigeonhole, there will be no such case.

 Case 2: In the bottom left 3x2 matrix of the 4x4 grid, two rows contain 1, -1 and one contains -1, 1.

 Sweet! Notice that when we have two rows that contain 1, -1, then the third row is uniquely determined as -1, 1. Thus, we just need to count the cases of 1,-1 and we got our answer.

 There are two -1's in the first column, and two 1's in the second. Thus, our orderings (in respective row order), must look like some permutation of
 \[\begin{bmatrix} -1 & 1 \\ -1 & 1 \\ 1 & -1 \end{bmatrix}\]

 This is an easy permutation! Say that A = \( \begin{bmatrix}-1 & 1\end{bmatrix} \) and B = \( \begin{bmatrix}1 & -1\end{bmatrix} \). Then, we have the matrix
 \[\begin{bmatrix} A \\ A \\ B \end{bmatrix}\]in which there are \(\frac{3!}{2!} = 3\) possible arrangements. Now, let's consider the bottom right 3x2 matrix of the 4x4 matrix in this case. For the first 3x1 matrix, we just need to distribute either two 1's or two -1's (doesn't matter, as it is just dependent on the top number), which can happen in \( 3 \choose 2 \) = 3 ways, and the last number is uniquely determined. Then, the last 3x1 matrix is uniquely determined, giving us a total of 3 \cdot 3 = 9 total cases in this case.

 Case 3: In the bottom left 3x2 matrix of the 4x4 grid, one row has -1, 1.

 In simpler terms, a case like
 \[\begin{bmatrix} -1 & 1 \\ 1 & 1 \\ -1 & -1 \end{bmatrix}\]satisfies the constraints.

 Label the cases A,B,C top down in respective order. Then, there are $3! = 6$ total configurations. In two of the rows (the rows containing B and C), they already determine the last two numbers of said rows. In the remaining row containing A, note that those numbers are also uniquely determined from the preceding rows, as they are already filled!

 However, notice that this case already covers our case 4, which one row contains 1,-1! This is due to the ordering of the top row (thank them later). If you aren't convinced, we can order the top row to make a case consisting of one row containing -1,1 possible, and we can reorder the top row again to make a case consisting of one row containing 1,-1. For example, we can arrange our top row as
 \[\begin{bmatrix} 1 & -1 & -1 & 1 \end{bmatrix}\]to obtain the 3x2 matrix
 \[\begin{bmatrix} -1 & 1 \\ 1 & 1 \\ -1 & -1 \end{bmatrix}\], but we can also order the top row as
 \[\begin{bmatrix} -1 & 1 & -1 & 1 \end{bmatrix}\]to obtain the matrix
 \[\begin{bmatrix} -1 & -1 \\ 1 & 1 \\ 1 & -1 \end{bmatrix}\]

 Boom! Two birds with one stone! (This is also called a 1-1 correspondence or bijection).

 Thus, there are 6 total cases for this case.

 We have a total of 6 possible orderings for the first row. Then, two cases emerge, in which, because each of the 6 possible orderings of the first row can affect the number of cases, we multiply the sum of each case by the total number of orderings of the first row. Thus, our answer is \( 6 \cdot (3 \cdot 3 + 6) = 6 \cdot (15) = \boxed{090} \).

 


	\item (\textit{AIME I 1997}) Given a nonnegative real number $x$, let $\langle x\rangle$ denote the fractional part of $x$; that is, $\langle x\rangle=x-\lfloor x\rfloor$, where $\lfloor x\rfloor$ denotes the greatest integer less than or equal to $x$. Suppose that $a$ is positive, $\langle a^{-1}\rangle=\langle a^2\rangle$, and $2<a^2<3$. Find the value of $a^{12}-144a^{-1}$.


 
 \textbf{Answer: }233



\textbf{Solution 1}

Looking at the properties of the number, it is immediately guess-able that $a = \phi = \frac{1+\sqrt{5}}2$ (the \href{https://artofproblemsolving.com/wiki/index.php?title=Phi}{golden ratio}) is the answer. The following is the way to derive that:

 Since $\sqrt{2} < a < \sqrt{3}$, $0 < \frac{1}{\sqrt{3}} < a^{-1} < \frac{1}{\sqrt{2}} < 1$. Thus $\langle a^2 \rangle = a^{-1}$, and it follows that $a^2 - 2 = a^{-1} \Longrightarrow a^3 - 2a - 1 = 0$. Noting that $-1$ is a root, this factors to $(a+1)(a^2 - a - 1) = 0$, so $a = \frac{1 + \sqrt{5}}{2}$ (we discard the negative root).

 Our answer is $(a^2)^{6}-144a^{-1} = \left(\frac{3+\sqrt{5}}2\right)^6 - 144\left(\frac{2}{1 + \sqrt{5}}\right)$. \href{https://artofproblemsolving.com/wiki/index.php?title=Complex\_conjugate}{Complex conjugates} reduce the second term to $-72(\sqrt{5}-1)$. The first term we can expand by the \href{https://artofproblemsolving.com/wiki/index.php?title=Binomial\_theorem}{binomial theorem} to get $\frac 1{2^6}\left(3^6 + 6\cdot 3^5\sqrt{5} + 15\cdot 3^4 \cdot 5 + 20\cdot 3^3 \cdot 5\sqrt{5} + 15 \cdot 3^2 \cdot 25 + 6 \cdot 3 \cdot 25\sqrt{5} + 5^3\right)$ $= \frac{1}{64}\left(10304 + 4608\sqrt{5}\right) = 161 + 72\sqrt{5}$. The answer is $161 + 72\sqrt{5} - 72\sqrt{5} + 72 = \boxed{233}$.

 Note that to determine our answer, we could have also used other properties of $\phi$ like $\phi^3 = 2\phi + 1$.

 



\textbf{Solution 2}

Find $a$ as shown above. Note that, since $a$ is a root of the equation $a^3 - 2a - 1 = 0$, $a^3 = 2a + 1$, and $a^{12} = (2a + 1)^4$. Also note that, since $a$ is a root of $a^2 - a - 1 = 0$, $\frac{1}{a} = a - 1$. The expression we wish to calculate then becomes $(2a + 1)^4 - 144(a - 1)$. Plugging in $a = \frac{1 + \sqrt{5}}{2}$, we plug in to get an answer of $(161 + 72\sqrt{5}) + 72 - 72\sqrt{5} = 161 + 72 = \boxed{233}$.

 



\textbf{Solution 3}

Find $a$ as shown above. Note that $a$ satisfies the equation $a^2 = a+1$ (this is the equation we solved to get it). Then, we can simplify $a^{12}$ as follows using the fibonacci numbers:

 $a^{12} = a^{11}+a^{10}= 2a^{10} + a^{9} = 3a^9+ 2a^8 = ... = 144a^1+89a^0 = 144a+89$

 So we want $144(a-\frac1a)+89 = 144(1)+89 = \boxed{233}$ since $a-\frac1a = 1$ is equivalent to $a^2 = a+1$.

 



\textbf{Solution 4}

As Solution 1 stated, $a^3 - 2a - 1 = 0$. $a^3 - 2a - 1 = a^3 - a^2 -a + a^2 -a -1 = (a+1)(a^2 - a - 1)$. So, $a^2 - a - 1 = 0$, $1 = a^2 - a$, $\frac1a = a-1$, $a^3 = 2a+1$, $a^2 = a+1$.

 $a^6 = (a^3)^2 = (2a+1)^2= 4a^2 + 4a +1= 4(a+1) + 4a + 1= 8a+5$
$a^{12} = (a^6)^2 = (8a+5)^2 = 64a^2 + 80a + 25 = 64(a+1) + 80a + 25 = 144a + 89$

 Therefore, $a^{12} - 144 a^{-1} = 144a + 89 - 144(a-1) = 89 + 144 = \boxed{233}$

 
Another way to factor $a^3 - 2a - 1$:

 $a^3 - 2a - 1 = a^3 + 1 -2a -2 = (a+1)(a^2 - a + 1) - 2(a+1) = (a + 1)(a^2 - a - 1)$




\textbf{Solution 5 (Dangerous)}

We can estimate $a$ to 3 decimal levels of precision by plugging it in to $a^3 = 3a + 1$ (insane bash) when we get a is slightly below $1.62$. We then use $a^3 = 3a + 1$ to manipulate what we want to get ( which is $88(a - a^2) + 321$ and by plugging $1.61$ and $1.62$ to in make sense for the answer to be $\boxed{233}$ rather than $234$

 Sidenote: I was really stupid and did not factor the cubic correctly because I thought it was not factorable. So do not bash it out this way.

 


	\item (\textit{AIME I 1997}) Every card in a deck has a picture of one shape - circle, square, or triangle, which is painted in one of the three colors - red, blue, or green. Furthermore, each color is applied in one of three shades - light, medium, or dark. The deck has 27 cards, with every shape-color-shade combination represented. A set of three cards from the deck is called complementary if all of the following statements are true:


 i. Either each of the three cards has a different shape or all three of the card have the same shape.


 ii. Either each of the three cards has a different color or all three of the cards have the same color.


 iii. Either each of the three cards has a different shade or all three of the cards have the same shade.


 How many different complementary three-card sets are there?


 
 \textbf{Answer: }117



\textbf{Solution 1}


            \begin{itemize}
        \item \textbf{Case 1}: All three attributes are the same. This is impossible since sets contain distinct cards.
        \item \textbf{Case 2}: Two of the three attributes are the same. There are ${3\choose 2}$ ways to pick the two attributes in question. Then there are $3$ ways to pick the value of the first attribute, $3$ ways to pick the value of the second attribute, and $1$ way to arrange the positions of the third attribute, giving us ${3\choose 2} \cdot 3 \cdot 3 = 27$ ways.
        \item \textbf{Case 3}: One of the three attributes are the same. There are ${3\choose 1}$ ways to pick the one attribute in question, and then $3$ ways to pick the value of that attribute. Then there are $3!$ ways to arrange the positions of the next two attributes, giving us ${3\choose 1} \cdot 3 \cdot 3! = 54$ ways.
        \item \textbf{Case 4}: None of the three attributes are the same. We fix the order of the first attribute, and then there are $3!$ ways to pick the ordering of the second attribute and $3!$ ways to pick the ordering of the third attribute. This gives us $(3!)^2 = 36$ ways.
            \end{itemize}
            Adding the cases up, we get $27 + 54 + 36 = \boxed{117}$.

 



\textbf{Solution 2}

Let's say we have picked two cards. We now compare their attributes to decide how we can pick the third card to make a complement set. For each of the three attributes, should the two values be the same we have one option - choose a card with the same value for that attribute. Furthermore, should the two be different there is only one option- choose the only value that is remaining. In this way, every two card pick corresponds to exactly one set, for a total of $\binom{27}{2} = 27*13 = 351$ possibilities. Note, however, that each set is generated by ${3\choose 2} = 3$ pairs, so we've overcounted by a multiple of 3 and the answer is $\frac{351}{3} = \boxed{117}$.

 



\textbf{Solution 3}

Treat the sets as ordered. Then for each of the three criterion, there are $3!=6$ choices if the attribute is different and there are $3$ choices is the attribute is the same. Thus all three attributes combine to a total of $(6+3)^3=729$ possibilities. However if all three attributes are the same then the set must be composed of three cards that are the same, which is impossible. This takes out $3^3=27$ possibilities. Notice that we have counted every set $3!=6$ times by treating the set as ordered. The final solution is then $\frac{729-27}{6}=\boxed{117}$




\textbf{Solution 4 (rewording of Solution 3)}

Note that for a possible valid triplet, two cards uniquely determine the third. $27$ choices for each card, so we have $27^2 = 729$ without any conditions. However, we have overcounted. The invalid triplets arise when all three cards are identical to each other, which can only happen in $3^3$ ways ($3$ ways to choose each of the attributes), so we have $729-27=702.$ However, we have overcounted by treating these as ordered, so we divide by $6,$ yielding $\frac{702}{6}=\boxed{117}.$$\newline$~martianrunner

 


	\item (\textit{AIME I 1997}) Let $x=\frac{\sum\limits_{n=1}^{44} \cos n^\circ}{\sum\limits_{n=1}^{44} \sin n^\circ}$. What is the greatest integer that does not exceed $100x$?


 
 \textbf{Answer: }241



\textbf{Solution 1}

Note that $\frac{\sum_{n=1}^{44} \cos n}{\sum_{n=1}^{44} \sin n} = \frac{\sum_{n=1}^{44} \cos n}{\sum_{n=46}^{89} \cos n} = \frac {\cos 1 + \cos 2 + \dots + \cos 44}{\cos 89 + \cos 88 + \dots + \cos 46}$ by the cofunction identities.(We could have also written it as $\frac{\sum_{n=1}^{44} \cos n}{\sum_{n=1}^{44} \sin n} = \frac{\sum_{n=46}^{89} \sin n}{\sum_{n=1}^{44} \sin n} = \frac {\sin 89 + \sin 88 + \dots + \sin 46}{\sin 1 + \sin 2 + \dots + \sin 44}$.)

 Now use the sum-product formula $\cos x + \cos y = 2\cos\left(\frac{x+y}{2}\right)\cos\left(\frac{x-y}{2}\right)$ We want to pair up $[1, 44]$, $[2, 43]$, $[3, 42]$, etc. from the numerator and $[46, 89]$, $[47, 88]$, $[48, 87]$ etc. from the denominator. Then we get:
 \[\frac{\sum_{n=1}^{44} \cos n}{\sum_{n=1}^{44} \sin n} = \frac{\sum_{n=1}^{44} \cos n}{\sum_{n=46}^{89} \cos n} = \frac{2\cos(\frac{45}{2})[\cos(\frac{43}{2})+\cos(\frac{41}{2})+\dots+\cos(\frac{1}{2})]}{2\cos(\frac{135}{2})[\cos(\frac{43}{2})+\cos(\frac{41}{2})+\dots+\cos(\frac{1}{2})]} \Rightarrow \frac{\cos(\frac{45}{2})}{\cos(\frac{135}{2})}\]

 To calculate this number, use the half angle formula. Since $\cos\left(\frac{x}{2}\right) = \pm \sqrt{\frac{\cos x + 1}{2}}$, then our number becomes: 
 \[\frac{\sqrt{\frac{\frac{\sqrt{2}}{2} + 1}{2}}}{\sqrt{\frac{\frac{-\sqrt{2}}{2} + 1}{2}}}\] in which we drop the negative roots (as it is clear cosine of 

 


	\item (\textit{AIME I 1997}) The function $f$ defined by $f(x)= \frac{ax+b}{cx+d}$, where $a$, $b$, $c$ and $d$ are nonzero real numbers, has the properties $f(19)=19$, $f(97)=97$ and $f(f(x))=x$ for all values except $\frac{-d}{c}$. Find the unique number that is not in the range of $f$.


 
 \textbf{Answer: }058



\textbf{Solution 1}

First, we use the fact that $f(f(x)) = x$ for all $x$ in the domain. Substituting the function definition, we have $\frac {a\frac {ax + b}{cx + d} + b}{c\frac {ax + b}{cx + d} + d} = x$, which reduces to
 \[\frac {(a^2 + bc)x + b(a + d)}{c(a + d)x + (bc + d^2)} =\frac {px + q}{rx + s} = x.\]In order for this fraction to reduce to $x$, we must have $q = r = 0$ and $p = s\not = 0$. From $c(a + d) = b(a + d) = 0$, we get $a = - d$ or $b = c = 0$. The second cannot be true, since we are given that $a,b,c,d$ are nonzero. This means $a = - d$, so $f(x) = \frac {ax + b}{cx - a}$.

 The only value that is not in the range of this function is $\frac {a}{c}$. To find $\frac {a}{c}$, we use the two values of the function given to us. We get $2(97)a + b = 97^2c$ and $2(19)a + b = 19^2c$. Subtracting the second equation from the first will eliminate $b$, and this results in $2(97 - 19)a = (97^2 - 19^2)c$, so
 \[\frac {a}{c} = \frac {(97 - 19)(97 + 19)}{2(97 - 19)} = 58 .\]

 Alternatively, we could have found out that $a = -d$ by using the fact that $f(f(-b/a))=-b/a$.

 



\textbf{Solution 2}

First, we note that $e = \frac ac$ is the horizontal \href{https://artofproblemsolving.com/wiki/index.php?title=Asymptote\_(Geometry)}{asymptote} of the function, and since this is a linear function over a linear function, the unique number not in the range of $f$ will be $e$. $\frac{ax+b}{cx+d} = \frac{b-\frac{cd}{a}}{cx+d} + \frac{a}{c}$. \href{https://artofproblemsolving.com/wiki/index.php?title=Without\_loss\_of\_generality}{Without loss of generality}, let $c=1$, so the function becomes $\frac{b- \frac{d}{a}}{x+d} + e$.  

 (Considering $\infty$ as a limit) By the given, $f(f(\infty)) = \infty$. $\lim_{x \rightarrow \infty} f(x) = e$, so $f(e) = \infty$. $f(x) \rightarrow \infty$ as $x$ reaches the vertical \href{https://artofproblemsolving.com/wiki/index.php?title=Asymptote\_(Geometry)}{asymptote}, which is at $-\frac{d}{c} = -d$. Hence $e = -d$. Substituting the givens, we get 

 
 \begin{align*} 19 &= \frac{b - \frac da}{19 - e} + e\\ 97 &= \frac{b - \frac da}{97 - e} + e\\ b - \frac da &= (19 - e)^2 = (97 - e)^2\\ 19 - e &= \pm (97 - e) \end{align*}


 Clearly we can discard the positive root, so $e = 58$.

 



\textbf{Solution 3}

We first note (as before) that the number not in the range of
 \[f(x) = \frac{ax+b}{cx+ d} = \frac{a}{c} + \frac{b - ad/c}{cx+d}\]is $a/c$, as $\frac{b-ad/c}{cx+d}$ is evidently never 0 (otherwise, $f$would be a constant function, violating the condition $f(19) \neq f(97)$).

 We may represent the real number $x/y$ as$\begin{pmatrix}x \\ y\end{pmatrix}$, with two such \href{https://artofproblemsolving.com/wiki/index.php?title=Vector}{column vectors}considered equivalent if they are scalar multiples of each other.  Similarly,we can represent a function $F(x) = \frac{Ax + B}{Cx + D}$ as a matrix$\begin{pmatrix} A & B\\ C& D \end{pmatrix}$.  Function composition andevaluation then become matrix multiplication.

 Now in general,
 \[f^{-1} = \begin{pmatrix} a & b\\ c&d \end{pmatrix}^{-1} = \frac{1}{\det(f)} \begin{pmatrix} d & -b \\ -c & a \end{pmatrix} .\]In our problem $f^2(x) = x$.  It follows that
 \[\begin{pmatrix} a & b \\ c& d \end{pmatrix} = K \begin{pmatrix} d & -b \\ -c & a \end{pmatrix} ,\]for some nonzero real $K$.  Since
 \[\frac{a}{d} = \frac{b}{-b} = K,\]it follows that $a = -d$.  (In fact, this condition condition is equivalentto the condition that $f(f(x)) = x$ for all $x$ in the domain of $f$.)

 We next note that the function
 \[g(x) =  x - f(x) = \frac{c x^2 + (d-a) x - b}{cx + d}\]evaluates to 0 when $x$ equals 19 and 97.  Therefore
 \[\frac{c x^2 + (d-a) x - b}{cx+d} = g(x) = \frac{c(x-19)(x-97)}{cx+d}.\]Thus $-19 - 97 = \frac{d-a}{c} = -\frac{2a}{c}$, so $a/c = (19+97)/2 = 58$,our answer.

 



\textbf{Solution 4}

Any number that is not in the domain of the inverse of $f(x)$ cannot be in the range of $f(x)$. Starting with $f(x) = \frac{ax+b}{cx+d}$, we rearrange some things to get $x = \frac{b-f(x)d}{f(x)c-a}$. Clearly, $\frac{a}{c}$ is the number that is outside the range of $f(x)$.

 
Since we are given $f(f(x))=x$, we have that 
 \[x = \frac{a\frac{ax+b}{cx+d}+b}{c\frac{ax+b}{cx+d}+d} = \frac{a^2x +ab+bcx+bd}{acx+bc+cdx+d^2} = \frac{x(bc+a^2)+b(a+d)}{cx(a+d)+(bc+d^2)}\]
 \[cx^2(a+d)+x(bc+d^2) = x(bc+a^2) + b(a+d)\]All the quadratic terms, linear terms, and constant terms must be equal on both sides for this to be a true statement so we have that $a = -d$.

 This solution follows in the same manner as the last paragraph of the first solution.

 



\textbf{Solution 5}

Since $f(f(x))$ is $x$, it must be symmetric across the line $y=x$. Also, since $f(19)=19$, it must touch the line $y=x$ at $(19,19)$ and $(97,97)$. $f$ a hyperbola that is a scaled and transformed version of $y=\frac{1}{x}$. Write $f(x)= \frac{ax+b}{cx+d}$ as $\frac{y}{cx+d}+z$, and z is our desired answer $\frac{a}{c}$. Take the basic hyperbola, $y=\frac{1}{x}$. The distance between points $(1,1)$ and $(-1,-1)$ is $2\sqrt{2}$, while the distance between $(19,19)$ and $(97,97)$ is $78\sqrt{2}$, so it is $y=\frac{1}{x}$ scaled by a factor of $39$. Then, we will need to shift it from $(-39,-39)$ to $(19,19)$, shifting up by $58$, or $z$, so our answer is $\boxed{58}$. Note that shifting the $x$ does not require any change from $z$; it changes the denominator of the part $\frac{1}{x-k}$.

 



\textbf{Solution 6 (Short)}

From $f(f(x))=x$, it is obvious that $\frac{-d}{c}$ is the value not in the range. First notice that since $f(0)=\frac{b}{d}$, $f(\frac{b}{d})=0$ which means $a(\frac{b}{d})+b=0$ so $a=-d$. Using $f(19)=19$, we have that $b=361c+38d$; on $f(97)=97$ we obtain $b=9409c+194d$. Solving for $d$ in terms of $c$ leads us to $d=-58c$, so the answer is $\boxed{058}$. 

 



\textbf{Solution 7}

Begin by finding the inverse function of $f(x)$, which turns out to be $f^{-1}(x)=\frac{19d-b}{a-19c}$. Since $f(f(x))=x$, $f(x)=f^{-1}(x)$, so substituting 19 and 97 yields the system, $\begin{array}{lcl} \frac{19a+b}{19c+d} & = & \frac{19d-b}{a-19c} \\ \frac{97a+b}{97c+d} & = & \frac{97d-b}{a-97c} \end{array}$, and after multiplying each equation out and subtracting equation 1 from 2, and after simplifying, you will get $116c=a-d$. Coincidentally, then $116c+d=a$, which is familiar because $f(116)=\frac{116a+b}{116c+d}$, and since $116c+d=a$, $f(116)=\frac{116a+b}{a}$. Also, $f(f(116))=\frac{a(\frac{116a+b}{a})+b}{c(\frac{116a+b}{a})+d}=116$, due to $f(f(x))=x$. This simplifies to $\frac{116a+2b}{c(\frac{116a+b}{a})+d}=116$, $116a+2b=116(c(\frac{116a+b}{a})+d)$, $116a+2b=116(c(116+\frac{b}{a})+d)$, $116a+2b=116c(116+\frac{b}{a})+116d$, and substituting $116c+d=a$ and simplifying, you get $2b=116c(\frac{b}{a})$, then $\frac{a}{c}=58$. Looking at $116c=a-d$ one more time, we get $116=\frac{a}{c}+\frac{-d}{c}$, and substituting, we get $\frac{-d}{c}=\boxed{58}$, and we are done.

 



\textbf{Solution 8 (shorter than solution 6)}

Because there are no other special numbers other than $19$ and $97$, take the average to get $\boxed{58}$. (Note I solved this problem the solution one way but noticed this and this probably generalizes to all $f(x)=x, f(y)=y$ questions like these)

 



\textbf{Solution 9 (Simple)}

By the function definition, $f(f(x))$, $f$ is its own inverse, so the only value not in the range of $f$ is the value not in the domain of $f$ (which is $-d/c$). 

 Since $f(f(x))$, $f(f(0)=0$ (0 is a convenient value to use). $f(f(0))=f(f(\frac{b}{d})=\dfrac{a\cdot\frac{b}{d}+b}{c\cdot\frac{b}{d}+d}=\dfrac{ab+bd}{bc+d^2}=0 \Rightarrow ab+bd=0$.

 Then$ab+bd=b(a+d)=0$ and since $b$ is nonzero, $a=-d$. 

 The answer we are searching for, $\dfrac{-d}{c}$ (the only value not in the range of $f$), can now be expressed as $\dfrac{a}{c}$.

 We are given $f(19)=19$ and $f(97)$, and they satisfy the equation $f(x)=x$, which simplifies to $\dfrac{ax+b}{cx+d}=x\Rightarrow x(cx+d)=ax+b\Rightarrow cx^2+(d-a)x+b=0$. We have written this quadratic with roots $19$ and $97$. 

 By Vieta, $\dfrac{-(d-a)}{c}=\dfrac{-(-a-a)}{c}=\dfrac{2a}{c}=19+97$.

 So our answer is $\dfrac{116}{2}=\boxed{058}$.

 



\textbf{Solution 9 (30-sec solve)}

Notice that the function is just an \href{https://artofproblemsolving.com/wiki/index.php?title=Involution}{involution} on the real number line. Since the involution has two fixed points, namely $19$ and $97$, we know that the involution is an \href{https://artofproblemsolving.com/wiki/index.php?title=Circular\_inversion}{inversion} with respect to a circle with a diameter from $19$ to $97$. The only point that is undefined under an inversion is the center of the circle, which we know is $\frac{19+97}{2}=\boxed{58}$ in both $x$ and $y$ dimensions.

 
Or if you don't think about inversion: A linear rational function like this is $f= a/c + (b-d)/(cx+d)$, and so has asymptotes at $x=-d/c$ and $y=a/c$, and these values must be equal because $f$ is an "involution", its own inverse. (Reflecting over $x=y$ does not change $f$). 

 By self-inverse symmetry, both asymptotes are equidistant to the graph points $(19,19)$ and $(97,97)$, so they must intersect at the mean of $19$ and $97$, which is $58$.

 



\textbf{Solution 10}

First, consider the equation $\frac{ax+b}{cx+d}=x$. This is a quadratic in terms of $x$, and we are given that $19,97$ are solutions to this equation. Rearranging yields $cx^2-(a-d)x-b=0$, so by Vieta’s Formulas we must have $\cfrac{a-d}{c}=19+97=116$.

 Next, we consider the second condition. Note that for $m=-\cfrac{b}{a},f(m)=0$, so $f(f(m))=f(0)$. We are also given that $f(f(m))=m$, so $f(0)=m$, which simplifies to $\cfrac{b}{d}=-\cfrac{b}{a}$ and $a=-d$. Substituting this into the equation above yields
 \[\cfrac{a+a}{c}=116\]
 \[\cfrac{2a}{c}=116\]
 \[\cfrac{a}{c}=58\]This is the horizontal asymptote of the function, so the function does not include this $y$-value. Thus the answer is $\boxed{58}$. ~eevee9406

 



\textbf{Solution 11 (Pure Computation)}

We first substitute in $19$ and $97$.

 We obtain

 
 \[\frac{19a+b}{19c+d}=19, \frac{97a+b}{97c+d}=97.\]

 Since $f(f(x))=x$, we know that 

 
 \[\frac{af(x)+b}{cf(x)+d} = x.\]

 Multiplying by $cf(x)+d$ gives $cxf(x)+dx=af(x)+b$, so $(a-cx)f(x)=dx-b$.

 Thus, $f(x)=\frac{dx-b}{a-cx}=\frac{ax+b}{cx+d}$.

 After cross multiplying, we obtain

 
 \[-acx^2+(a^2-bc)x+ab=cdx^2+(d^2-bc)x-bd.\]

 Since this holds for all $x$, the coefficients are equal. Thus, $-ac=cd$, $a^2-bc=d^2-bc$, and $ab=-bd$.

 This gives $-a=d$.

 Substituting this into $\frac{19a+b}{19c+d}=19$ gives

 
 \[361c-19a=19a+b \implies 38a=361c-b.\]

 Applying the same logic to $\frac{97a+b}{97c+d}=97$, $194a=9409c-b$.

 Subtracting the first equation from the second gives

 
 \[156a = 9048c \implies a = 58c.\]

 Since $a=58c$, then $d=-58c$.

 From $38a=361c-b$, we get $b=361c-38a$. Since $a=58c$, simplifying gives $b=-1843c$.

 Thus,

 
 \[f(x)=\frac{58cx-1843c}{cx-58c} \implies f(x)=\frac{58x-1843}{x-58}.\]

 Note that this step is valid because $c$ is nonzero.

 Then, we can split this fraction into

 
 \[58+\frac{1521}{x-58}.\]

 Since a defined fraction is equal to $0$ if and only if its numerator is $0$, we know that $\frac{1521}{x-58}$ is nonzero.

 Thus, $58$ cannot be obtained from $f(x)$.

 Finally, our answer is $\boxed{058}$.

 


	\item (\textit{AIME I 1997}) 
Let $S$ be the set of points in the Cartesian plane that satisfy 

 $\Big|\big| |x|-2\big|-1\Big|+\Big|\big| |y|-2\big|-1\Big|=1.$

  If a model of $S$ were built from wire of negligible thickness, then the total length of wire required would be $a\sqrt{b}$, where $a$ and $b$ are positive integers and $b$ is not divisible by the square of any prime number. Find $a+b$.


 
 \textbf{Answer: }066



\textbf{Solution 1 (non-rigorous)}

Let $f(x) = \Big|\big||x|-2\big|-1\Big|$, $f(x) \ge 0$. Then $f(x) + f(y) = 1 \Longrightarrow f(x), f(y) \le 1 \Longrightarrow x, y \le 4$. We only have a $4\times 4$ area, so guessing points and graphing won't be too bad of an idea. Since $f(x) = f(-x)$, there's a symmetry about all four \href{https://artofproblemsolving.com/wiki/index.php?title=Quadrant}{quadrants}, so just consider the first quadrant. We now gather some points:

 


	\item (\textit{AIME I 1997}) Let $v$ and $w$ be distinct, randomly chosen roots of the equation $z^{1997}-1=0$. Let $m/n$ be the probability that $\sqrt{2+\sqrt{3}}\le |v+w|$, where $m$ and $n$ are relatively prime positive integers. Find $m+n$.


 
 \textbf{Answer: }582



\textbf{Solution 1}


            \begin{enumerate}[label=(\roman*)]
                \item $z^{1997}=1=1(\cos 0 + i \sin 0)$
            \end{enumerate}
            Define $\theta = 2\pi/1997$. By \href{https://artofproblemsolving.com/wiki/index.php?title=De\_Moivre\%27s\_Theorem}{De Moivre's Theorem} the roots are given by

 
            \begin{enumerate}[label=(\roman*)]
                \item $z=\cos (k\theta) +i\sin(k\theta), \qquad k \in \{0,1,\ldots,1996\}$
            \end{enumerate}
            Now, let $v$ be the root corresponding to $m\theta=2m\pi/1997$, and let $w$ be the root corresponding to $n\theta=2n\pi/ 1997$. Then 
 \begin{align*} |v+w|^2 &= \left(\cos(m\theta) + \cos(n\theta)\right)^2 + \left(\sin(m\theta) + \sin(n\theta)\right)^2 \\ &= 2 + 2\cos\left(m\theta\right)\cos\left(n\theta\right) + 2\sin\left(m\theta\right)\sin\left(n\theta\right) \end{align*}
The \href{https://artofproblemsolving.com/wiki/index.php?title=Trigonometric\_identities}{cosine difference identity} simplifies that to
 \[|v+w|^2 = 2+2\cos((m-n)\theta)\]

 We need $|v+w|^2 \ge 2+\sqrt{3}$, which simplifies to 
 \[\cos((m-n)\theta) \ge \frac{\sqrt{3}}{2}.\] Hence, $-\frac{\pi}{6} \le (m-n) \theta \le \frac{\pi}{6},$ or alternatively, $|m - n| \theta \le \frac{\pi}{6}.$ 

 Recall that $\theta=\frac{2\pi}{1997}.$ Thus, 
 \[|m - n| \le \frac{\pi}{6} \cdot \frac{1997}{2 \pi} = \left\lfloor \frac{1997}{12} \right\rfloor =166.\]

 Therefore, $m$ and $n$ cannot be more than $166$ away from each other.  This means that for a given value of $m$, there are $332$ values for $n$ that satisfy the inequality; $166$ of them $> m$, and $166$ of them $< m$ ($m \neq n$ since the problem says that $m$ and $n$ are distinct).  Since $m$ and $n$ must be distinct, $n$ can have $1996$ total possible values for each value of $m$.  Therefore, the probability is $\frac{332}{1996}=\frac{83}{499}$.  The answer is then $499+83=\boxed{582}$.

 



\textbf{Solution 2}

The solutions of the equation $z^{1997} = 1$ are the $1997$th \href{https://artofproblemsolving.com/wiki/index.php?title=Roots\_of\_unity}{roots of unity} and are equal to $\text{cis}(\theta_k)$, where $\theta_k = \frac {2\pi k}{1997}$ for $k = 0,1,\ldots,1996.$ Thus, they are located at uniform intervals on the unit circle in the complex plane. 

 The quantity $|v+w|$ is unchanged upon rotation around the origin, so, WLOG, we can assume $v=1$ after rotating the axis till $v$ lies on the real axis. Let $w=\text{cis}(\theta_k)$. Since $w\cdot \overline{w}=|w|^2=1$ and $w+\overline{w}=2\text{Re}(w) = 2\cos\theta_k$, we have 
 \[|v + w|^2  =  (1+w)(1+\overline{w}) = 2+2\cos\theta_k\]We want $|v + w|^2\ge 2 + \sqrt {3}.$  From what we just obtained, this is equivalent to 
 \[\cos\theta_k\ge \frac {\sqrt {3}}2 \qquad \Leftrightarrow \qquad -\frac {\pi}6\le \theta_k \le  \frac {\pi}6\] which is satisfied by $k = 166,165,\ldots, - 165, - 166$ (we don't include 0 because that corresponds to $v$).  So out of the $1996$ possible $k$, $332$ work.  Thus, $m/n = 332/1996 = 83/499.$  So our answer is $83 + 499 = \boxed{582}.$




\textbf{Solution 3}

We can solve a geometrical interpretation of this problem.

 Without loss of generality, let $w = 1$. We are now looking for a point exactly one unit away from $w$ such that the point is at least $\sqrt{2 + \sqrt{3}}$ units away from the origin. Note that the "boundary" condition is when the point will be exactly $\sqrt{2+\sqrt{3}}$ units away from the origin; these points will be the intersections of the circle centered at $(1,0)$ with radius $1$ and the circle centered at $(0,0)$ with radius $\sqrt{2+\sqrt{3}}$. The equations of these circles are $(x-1)^2 + y^2 = 1$ and $x^2 + y^2 = 2 + \sqrt{3}$. Solving for $x$ yields $x = 1 + \frac{\sqrt{3}}{2}$. Clearly, this means that the real part of $v$ is greater than $\frac{\sqrt{3}}{2}$. Solving, we note that $332$ possible $v$s exist, meaning that $\frac{m}{n} = \frac{332}{1996} = \frac{83}{499}$. Therefore, the answer is $83 + 499 = \boxed{582}$.

 



\textbf{Solution 4}

Since $z^{1997}=1$, the roots will have magnitude $1$. Thus, the roots can be written as $\cos(\theta)+i\sin(\theta)$ and $\cos(\omega)+i\sin(\omega)$ for some angles $\theta$ and $\omega$. We rewrite the requirement as $\sqrt{2+\sqrt3}\le|\cos(\theta)+\cos(\omega)+i\sin(\theta)+i\sin(\omega)|$, which can now be easily manipulated to $2+\sqrt{3}\le(\cos(\theta)+\cos(\omega))^2+(\sin(\theta)+\sin(\omega))^2$. 

 WLOG, let $\theta = 0$. Thus, our inequality becomes $2+\sqrt{3}\le(1+\cos(\omega))^2+(\sin(\omega))^2$, $2+\sqrt{3}\le2+2\cos(\omega)$, and finally $\cos(\omega)\ge\frac{\sqrt{3}}{2}$. Obviously, $\cos(\frac{\pi}{6})=\frac{\sqrt{3}}{2}$, and thus it follows that, on either side of a given point, $\frac{1997}{12}\approx166$ points will work. The probability is $\frac{166\times2}{1996} = \frac{83}{499}$, and thus our requested sum is $\boxed{582}$~SigmaPiE

 



\textbf{Solution 5}

Let \(v = \operatorname{cis}\left(\frac{2\pi k}{1997}\right)\) and  \(w = \operatorname{cis}\left(\frac{2\pi c}{1997}\right)\), where \(\operatorname{cis}\left(\theta\right) = \cos(\theta) + i\sin(\theta)\), \(v \neq w\), and \(k,c \in \{0,1,2,\dots,1996\}\)

 We are given
 \[\sqrt{2+\sqrt{3}} \le |v+w|.\]Squaring both sides gives
 \[2+\sqrt{3} \le |v+w|^2.\]

 Expanding \(v+w\), we obtain
 \[v+w = \cos\!\left(\frac{2\pi k}{1997}\right) + \cos\!\left(\frac{2\pi c}{1997}\right) + i\!\left(\sin\!\left(\frac{2\pi k}{1997}\right) + \sin\!\left(\frac{2\pi c}{1997}\right)\right).\]

 Using sum-to-product identities,
 \[v+w = 2\cos\!\left(\frac{2\pi(k+c)}{1997}\right)\cos\!\left(\frac{2\pi(k-c)}{1997}\right) + i\!\left(2\sin\!\left(\frac{2\pi(k+c)}{1997}\right)\cos\!\left(\frac{2\pi(k-c)}{1997}\right)\right).\]

 Thus,
 \[|v+w|^2 = 4\cos^2\!\left(\frac{2\pi(k-c)}{1997}\right) \left(\cos^2\!\left(\frac{2\pi(k+c)}{1997}\right) + \sin^2\!\left(\frac{2\pi(k+c)}{1997}\right)\right).\]

 Since \(\sin^2\theta+\cos^2\theta=1\), we have
 \[|v+w|^2 = 4\cos^2\!\left(\frac{2\pi(k-c)}{1997}\right).\]

 Substituting into the inequality gives
 \[2+\sqrt{3} \le 4\cos^2\!\left(\frac{2\pi(k-c)}{1997}\right).\]

 Taking square roots,
 \[\sqrt{2+\sqrt{3}} \le 2\cos\!\left(\frac{2\pi(k-c)}{1997}\right) \quad \text{or} \quad-\sqrt{2+\sqrt{3}} \ge 2\cos\!\left(\frac{2\pi(k-c)}{1997}\right).\]

 Dividing by \(2\),
 \[\sqrt{\frac{2+\sqrt{3}}{4}} \le \cos\!\left(\frac{2\pi(k-c)}{1997}\right) \quad \text{or} \quad  -\sqrt{\frac{2+\sqrt{3}}{4}} \ge \cos\!\left(\frac{2\pi(k-c)}{1997}\right).\]

 Noting that
 \[\sqrt{\frac{2+\sqrt{3}}{4}} = \cos \frac{\pi}{12},\]we conclude that
 \[-\frac{\pi}{12} \le \frac{2\pi(k-c)}{1997} \le \frac{\pi}{12} \quad \text{or} \quad \frac{11\pi}{12} \le \frac{2\pi(k-c)}{1997} \le \frac{13\pi}{12}.\]

 These two cases are symmetric, so we consider the first and double the result.  Simplifying,
 \[-\frac{1997}{24} \le k-c \le \frac{1997}{24}.\]

 Since \(k-c\) is an integer,
 \[-83 \le k-c \le 83.\]

 Because \(k \neq c\), for each fixed \(k\) there are exactly \(166\) valid values of \(c\).This works modulo \(1997\) since the arguments lie on a circle (e.g. \(\cos(\frac{\pi}{12})=\cos(\frac{-23\pi}{12})\)).

 Thus, the probability for one case is \(\frac{166}{1996}\), and doubling gives \(\frac{332}{1996} = \frac{83}{499}.\) The problem asks for \(m+n\), so \(83+499=\boxed{582}.\)

 


	\item (\textit{AIME I 1997}) The sides of rectangle $ABCD$ have lengths $10$ and $11$. An equilateral triangle is drawn so that no point of the triangle lies outside $ABCD$. The maximum possible area of such a triangle can be written in the form $p\sqrt{q}-r$, where $p$, $q$, and $r$ are positive integers, and $q$ is not divisible by the square of any prime number. Find $p+q+r$.


 
 \textbf{Answer: }554



\textbf{Solution 1 (Coordinate Bash)}

Consider points on the \href{https://artofproblemsolving.com/wiki/index.php?title=Complex\_plane}{complex plane} $A (0,0),\ B (11,0),\ C (11,10),\ D (0,10)$. Since the rectangle is quite close to a square, we figure that the area of the equilateral triangle is maximized when a vertex of the triangle coincides with that of the rectangle. Set one \href{https://artofproblemsolving.com/wiki/index.php?title=Vertex}{vertex} of the triangle at $A$, and the other two points $E$ and $F$ on $BC$ and $CD$, respectively. Let $E (11,a)$ and $F (b, 10)$. Since it's equilateral, then $E\cdot\text{cis}60^{\circ} = F$, so $(11 + ai)\left(\frac {1}{2} + \frac {\sqrt {3}}{2}i\right) = b + 10i$, and expanding we get $\left(\frac {11}{2} - \frac {a\sqrt {3}}{2}\right) + \left(\frac {11\sqrt {3}}{2} + \frac {a}{2}\right)i = b + 10i$.

 
                    \begin{figure}[H]
                        \centering
                        \includegraphics[width=0.5\textwidth]{Solution/AIME_I_1997_15_Solution_Figure_1.PNG}
                    \end{figure}
                    We can then set the real and imaginary parts equal, and solve for $(a,b) = (20 - 11\sqrt {3},22 - 10\sqrt {3})$. Hence a side $s$ of the equilateral triangle can be found by $s^2 = AE^2 = a^2 + AB^2 = 884 - 440\sqrt{3}$. Using the area formula $\frac{s^2\sqrt{3}}{4}$, the area of the equilateral triangle is $\frac{(884-440\sqrt{3})\sqrt{3}}{4} = 221\sqrt{3} - 330$. Thus $p + q + r = 221 + 3 + 330 = \boxed{554}$.

 



\textbf{Solution 2}

This is a trigonometric re-statement of the above. Let $x = \angle EAB$; by alternate interior angles, $\angle DFA=60+x$. Let $a = EB$ and the side of the equilateral triangle be $s$, so $s= \sqrt{a^2+121}$ by the \href{https://artofproblemsolving.com/wiki/index.php?title=Pythagorean\_Theorem}{Pythagorean Theorem}. Now $\frac{10}{s} = \sin(60+x)=  \sin {60} \cos x+ \cos {60} \sin x = \left(\frac{\sqrt{3}}2\right)\left(\frac{11}s\right)+\left(\frac 12\right)\left( \frac as \right)$. This reduces to $a=20-11\sqrt{3}$. 

 Thus, the area of the triangle is $\frac{s^2\sqrt{3}}{4} =(a^2+121)\frac{\sqrt{3}}{4}$, which yields the same answer as above.

 



\textbf{Solution 3}

Since $\angle{BAD}=90$ and $\angle{EAF}=60$, it follows that $\angle{DAF}+\angle{BAE}=90-60=30$. Rotate triangle $ADF$ $60$ degrees clockwise. Note that the image of $AF$ is $AE$. Let the image of $D$ be $D'$. Since angles are preserved under rotation, $\angle{DAF}=\angle{D'AE}$. It follows that $\angle{D'AE}+\angle{BAE}=\angle{D'AB}=30$. Since $\angle{ADF}=\angle{ABE}=90$, it follows that quadrilateral $ABED'$ is cyclic with circumdiameter $AE=s$ and thus circumradius $\frac{s}{2}$. Let $O$ be its circumcenter. By Inscribed Angles, $\angle{BOD'}=2\angle{BAD'}=60$. By the definition of circle, $OB=OD'$. It follows that triangle $OBD'$ is equilateral. Therefore, $BD'=r=\frac{s}{2}$. Applying the Law of Cosines to triangle $ABD'$, $\frac{s}{2}=\sqrt{10^2+11^2-(2)(10)(11)(\cos{30})}$. Squaring and multiplying by $\sqrt{3}$ yields $\frac{s^2\sqrt{3}}{4}=221\sqrt{3}-330\implies{p+q+r=221+3+330=\boxed{554}}$

 -Solution by \textbf{thecmd999}

 



\textbf{Solution 4 (Fast, no trig)}

Clearly one vertex of the equilateral triangle is on a vertex of the rectangle, and the other two are lying on two other sides. Let $m$ be the side length of the triangle, and let the rectangle be partitioned into the equilateral triangle, a right triangle with sides 11, $y$, $m$, a right triangle with sides 10, $x$, $m$, and a right triangle with sides $11-x$, $10-y$, $m$. Simple area analysis nets 
 \[110=\frac{\sqrt{3}}{4}m^2+\frac{11}{2}y+\frac{10}{2}x+\frac{(11-x)(10-y)}{2}\implies 110-\frac{\sqrt{3}}{2}m^2=xy\]By the Pythagorean Theorem, $11^2+y^2=m^2$ and $10^2+x^2=m^2$, so $x^2y^2=(m^2-10^2)(m^2-11^2)$. Thus,
 \[(110-\frac{\sqrt{3}}{2}m^2)^2=(m^2-10^2)(m^2-11^2)\]
 \[110^2-110\sqrt{3}m^2+\frac34m^4=m^4-221m^2+110^2\]Obviously $m^2\neq0$ so we can divide by $m^2$ after cancellation:
 \[-110\sqrt{3}+221=\frac14m^2\]The area of the triangle is $\frac{\sqrt{3}}{4}m^2$, so the finish is simple.
 \[\frac{\sqrt{3}}{4}m^2=221\sqrt{3}-330\implies p+q+r=221+3+330=\boxed{554}\]

 - clarkculus

 


\end{enumerate}